% ==============================================================================
% Section II: State of the Art
% ==============================================================================
\section{State of the Art}
\label{sec:related_work}

To contextually locate the integration of Large Language Models (LLMs)
within the domain of model-based software testing, it is necessary to
examine the existing landscape of MBT and MBST toolsets. Over several
decades, academic and commercial entities have developed specialized
testing engines designed to consume structured inputs and generate test
suites based on distinct coverage criteria~\cite{taxonomy_mbt,
mbt_tools_approach}. Table~\ref{tab:mbt_tools} summarises the most
prominent tools with respect to their input formalisms, generation
strategies, and target domains.

% ---------- Table I: MBT / MBST Tool Landscape ----------
\begin{table*}[!t]
\centering
\caption{Comparison of Established Model-Based Testing Tools}
\label{tab:mbt_tools}
\renewcommand{\arraystretch}{1.15}
\begin{tabular}{p{2.0cm} p{3.8cm} p{4.8cm} p{4.8cm}}
\hline\hline
\textbf{Tool} &
\textbf{Primary Input Format} &
\textbf{Test Generation Strategy \& Heuristics} &
\textbf{Target Domain \& Supported Outputs} \\
\hline
MaTeLo &
Markov chains, annotated sequence diagrams &
Profile-oriented random generation, transition coverage &
Functional test generation, TTCN-3, commercial systems \\
\hline
JUMBL &
Custom Test Model Language (TML) &
Automated constraint-driven probability generation &
Java command-line suite for statistical usage metrics \\
\hline
fMBT &
AAL pre/postcondition language &
Random, weighted random, and lookahead heuristics &
Open-source API and system testing with Python adapters \\
\hline
GraphWalker &
Finite State Machines (FSM) &
A* search, random walks with state/edge limits &
Open-source state machine path execution \\
\hline
Modbat &
Scala-based Domain-Specific Language (DSL) &
Probabilistic and non-deterministic transitions &
Specialised API testing, exception-handling verification \\
\hline
MoMuT::UML &
UML State Machines, Object-Oriented Action Systems &
Mutation-based, fault-driven test case generation &
Academic safety-critical and contract-based systems \\
\hline
BPM-X &
BPMN, UML business process models &
Statement, branch, path, and condition criteria &
Enterprise business process verification, ALM exports \\
\hline
Conformiq &
UML State Machines, QML, Activity Diagrams &
Combinatorial, constraint-driven test case generation &
Enterprise automated testing, test management platforms \\
\hline
4Test &
Gherkin-inspired custom syntax &
Constraint-driven combinatorial scenario selection &
Commercial interface and functional testing \\
\hline\hline
\end{tabular}
\end{table*}

A common thread across the tools listed in Table~\ref{tab:mbt_tools} is
their reliance on handcrafted, domain-specific input artefacts---Markov
chains~\cite{prowell_markov}, custom DSLs~\cite{jumbl_mbst_tutorial},
AAL annotations~\cite{boehr_constraints}, or full UML state
machines~\cite{mbt_overview_bme}. Constructing and maintaining these
artefacts remains the principal bottleneck for wider industrial
adoption~\cite{taxonomy_mbt}. It is precisely this bottleneck that
motivates the exploration of LLMs as a means to automate or partially
automate the generation of behavioural models from less formal sources.

% ---------- 2.1 Foundation Models for Formal Process Automation ----------
\subsection{Foundation Models for Formal Process Automation}
\label{sec:rw_foundation_models}

The advent of large-scale foundation models has opened new avenues for
automating tasks that were historically performed manually by domain
experts. Within protocol engineering, \emph{ProtocolGPT}~\cite{protocol_gpt}
demonstrated that Retrieval-Augmented Generation (RAG) can be used to
infer protocol state machines directly from RFC documents, achieving a
precision exceeding 90\% on well-structured specifications. By indexing
protocol documentation into a vector store and prompting the model with
targeted queries, ProtocolGPT effectively bridges the gap between
unstructured natural-language specifications and formal finite-state
representations.

Complementing this line of work, the \emph{ChatFuMe}
framework~\cite{chatfume} introduces an active-feedback parsing loop in
which the LLM iteratively refines its interpretation of protocol
documentation. Rather than relying on a single-pass generation,
ChatFuMe solicits structured clarifications and re-queries the model
until convergence, yielding more robust and consistent state-machine
extractions.

Beyond textual inputs, recent advances in multimodal LLMs have enabled
the interpretation of \emph{visual} specifications---hand-drawn
diagrams, whiteboard sketches, and scanned documents---into structured
UML artefacts~\cite{uml_benchmarking}. These capabilities suggest that
the input barrier for model-based testing could be further lowered by
accepting informal, heterogeneous documentation as a starting point for
model construction.

% ---------- 2.2 GUI Exploration and Visual Understanding ----------
\subsection{GUI Exploration and Visual Understanding}
\label{sec:rw_gui}

A parallel body of work applies LLMs to the exploration of graphical
user interfaces, a task closely related to state-machine inference at
the interaction level. \emph{VisionDroid}~\cite{visiondroid} employs
vision-language models to navigate mobile application screens and
automatically derive behavioural paths, while \emph{LLMShot} extends
this paradigm to few-shot GUI exploration with minimal human-provided
demonstrations. On the commercial side, Eggplant Functional's
\emph{Find By Description} engine~\cite{eggplant_ai} leverages natural
language to locate and interact with on-screen elements, effectively
enabling AI-driven exploratory testing without explicit object
repositories. These systems illustrate the growing convergence between
visual understanding, natural-language processing, and automated test
generation.

% ---------- TikZ Figure: LLM Pipeline ----------
\begin{figure}[!t]
\centering
\begin{tikzpicture}[
    node distance=1.6cm and 1.0cm,
    >=latex,
    every node/.style={font=\small},
    box/.style={draw, rounded corners=3pt, minimum height=0.9cm,
                minimum width=3.0cm, align=center, fill=white},
    llmbox/.style={draw, rounded corners=3pt, minimum height=3.4cm,
                   minimum width=5.6cm, align=center,
                   fill=gray!8, inner sep=8pt},
    resultbox/.style={draw, rounded corners=2pt, minimum height=0.7cm,
                      minimum width=4.8cm, align=center, fill=white,
                      font=\footnotesize},
]

% --- Input ---
\node[box, fill=blue!10] (prompt) {Natural Language\\Prompt};

% --- LLM container ---
\node[llmbox, right=1.8cm of prompt] (llm) {};
\node[anchor=north, font=\small\bfseries] at ([yshift=-4pt]llm.north)
    {Large Language Model};

% --- Two internal result boxes ---
\node[resultbox, fill=green!8] (single)
    at ([yshift=-2pt]llm.center |- {$(llm.north)!0.42!(llm.south)$})
    {Single-Prompt Baseline\\
     $F_1^{\mathrm{States}} \approx 0.80$\,,\;
     $F_1^{\mathrm{Trans}} \approx 0.75$};

\node[resultbox, fill=orange!10] (multi)
    at ([yshift=2pt]llm.center |- {$(llm.north)!0.75!(llm.south)$})
    {Multi-Step / Hybrid Framework\\
     $F_1^{\mathrm{States}} \approx 0.85$\,,\;
     $F_1^{\mathrm{Trans}} \approx 0.65$};

% --- Arrow ---
\draw[->, thick] (prompt.east) -- (llm.west);

\end{tikzpicture}
\caption{High-level LLM pipeline for state-machine extraction,
contrasting single-prompt and multi-step generation strategies with
representative $F_1$ scores for state and transition identification.}
\label{fig:llm_pipeline}
\end{figure}

% ---------- 2.3 Frameworks for LLM-Driven State Machine Generation --------
\subsection{Frameworks for LLM-Driven State Machine Generation}
\label{sec:rw_frameworks}

Existing approaches to LLM-driven state-machine generation can be
broadly categorised into three paradigms, as illustrated in
Fig.~\ref{fig:llm_pipeline}.

\subsubsection{Structure-Driven Frameworks}
Structure-driven methods impose an explicit schema or template on the
LLM's output. The model is instructed to populate predefined fields
(e.g., a JSON schema enumerating states, transitions, guards, and
actions), and the resulting artefact is validated against the schema
before downstream consumption~\cite{llm_state_machine_frameworks}. This
strategy favours \emph{syntactic correctness} and enables deterministic
post-processing, but may constrain the model's ability to capture
nuanced behavioural semantics that fall outside the template.

\subsubsection{Event-Driven Frameworks}
Event-driven approaches frame state-machine extraction as an
event-identification task: the LLM first enumerates the events or stimuli
described in the specification, and subsequently derives the states and
transitions that arise from those events~\cite{boundary_value_llm}. By
anchoring the generation process in observable events, these frameworks
align naturally with usage-model methodologies~\cite{usage_model_web,
nap_statistics} and can leverage the LLM's strength in entity extraction.

\subsubsection{Hybrid Approaches}
Hybrid strategies combine the immediacy of single-prompt generation with
structured, step-by-step refinement. A typical pipeline issues an
initial broad prompt to obtain a draft state machine, then applies a
sequence of targeted follow-up prompts---e.g., ``identify missing guard
conditions'' or ``merge equivalent states''---to iteratively improve the
model~\cite{llm_state_machine_frameworks}. As shown in
Fig.~\ref{fig:llm_pipeline}, hybrid approaches tend to achieve higher
$F_1$ scores for state identification (${\approx}\,0.85$) compared with
single-prompt baselines (${\approx}\,0.80$), albeit with a trade-off in
transition accuracy ($F_1^{\mathrm{Trans}} \approx 0.65$ vs.\
${\approx}\,0.75$), likely due to the compounding of intermediate
extraction errors across multiple inference steps.

\medskip
Collectively, these related works establish that LLMs possess a
demonstrable capacity to parse semi-structured and unstructured
specifications into formal behavioural models. However, a comprehensive,
end-to-end evaluation that benchmarks multiple prompting strategies
against an established MBST toolchain---and that measures not only model
fidelity but also downstream test-suite quality---remains an open gap in
the literature. The present study addresses this gap directly.
